\section{Обоснование выбранного метода}
\headingtextsep

В качестве основного метода развёртывания для системы «Контур» выбран контейнерный подход с использованием Docker Compose. Такое решение напрямую следует из архитектуры проекта и уже зафиксировано в исходных артефактах репозитория. В текущей реализации система включает PostgreSQL, API-сервис на Rust, web-клиент на React и отдельный MCP-мост. Все эти компоненты должны запускаться совместно и в согласованном порядке.

Выбор Docker Compose обоснован следующими причинами.

\begin{itemize}
\item Контейнерный подход обеспечивает единообразие среды запуска на разных машинах и снижает влияние различий в локальной конфигурации операционной системы.
\item Развёртывание многокомпонентной системы описывается в одном compose-файле, что делает процесс запуска прозрачным и повторяемым.
\item Каждый сервис получает собственный образ и собственный набор зависимостей, поэтому изменения в frontend, backend или интеграционном мосте меньше влияют друг на друга.
\item Наличие healthcheck для PostgreSQL и зависимостей между сервисами позволяет задать корректный порядок старта контейнеров.
\item Для первой версии проекта этот способ даёт хороший баланс между простотой эксплуатации и возможностью дальнейшего роста архитектуры.
\end{itemize}

Дополнительным аргументом является то, что выбранный метод уже соответствует документации и структуре репозитория. В проекте предусмотрены Dockerfile для основных компонентов и файл docker-compose.yml, который описывает запуск PostgreSQL, API, web-интерфейса и MCP-моста. Следовательно, именно контейнерный способ развёртывания является не только теоретически предпочтительным, но и практически поддерживаемым текущим состоянием системы.

Использование установщиков или пакетных менеджеров в данном случае привело бы к ручной координации нескольких процессов, росту числа зависимостей на машине и более высокой вероятности ошибок конфигурации. Поэтому для учебной, локальной и базовой self-hosted поставки выбор Docker Compose является наиболее рациональным.
